\chapter{UFO CSI}\label{ch:ufocsi}
\markboth{UFO CSI}{UFO CSI}

\chapquote{``Unlike scientology, science in the true sense of the word is open to unbiased investigation of any existing phenomena.''}{Stanislav Grof}

Everything in this chapter --- the taxonomies, the interview protocol, the instrument lists, the
insistence on baselines and denominators --- comes from one man, and if you learn nothing else
from these pages, learn him. Josef Allen Hynek was a professional astronomer, chairman of the
astronomy department at Northwestern, a man who spent his working life on stellar spectra and
satellite tracking. He is here because he is the only person in the history of this field who
turned an anecdote pile into a procedure.

I am deliberately not going to rehearse his government work here. That story --- what the Air
Force asked him to do, what it did with what he gave it, and how the whole arrangement ended ---
belongs to Chapter~\ref{ch:government}, and you should read it there. What matters for a chapter
about method is narrower and more useful: over roughly two decades of reading other people's
case files, Hynek worked out what a UFO report actually \emph{is} as an object of study, and
nobody before him had.

His answer was uncomfortable to both camps. \hi{A UFO report is not evidence about an object. It
is evidence about an observation} --- a witness, a set of viewing conditions, an instrument or
two if you are lucky, and a stimulus of unknown identity somewhere at the far end of the chain.
Treat it as a claim about a craft and you will argue forever. Treat it as a measurement with
known error sources and you can begin doing something. That single reframing is the whole basis
of civilian UFO investigation, and it is why this chapter exists.

From it he derived four rules that he stated, in one form or another, for the rest of his life.

\textbf{One: classify by observing conditions, not by excitement.} Before you sort reports by
what they might be, sort them by how they were seen --- distance, duration, lighting, number of
witnesses, instrumental corroboration. A brilliantly exciting account seen for two seconds by one
person is a worse data point than a dull one seen for four minutes by six people and a radar
scope, and any scheme that does not encode that fact will mislead you.

\textbf{Two: the residue is the phenomenon.} Most reports resolve. Hynek said so bluntly, and he
said it far more often than the people who quote him admit. What interested him was the
fraction that survived competent identification --- the residue --- and his argument was never
``therefore spacecraft''. It was that the residue existed, that its size was measurable, and that
a science which refuses to measure a residue it finds embarrassing has stopped being a science.

\textbf{Three: an explanation offered before an investigation is not an explanation.} He learned
this one the hard way and in public, and I will come back to it under interviewing technique. His
formulation for the general disease was the \hi{``will to disbelieve''}: the reflex that reaches
for any available answer, however poorly fitted, because the alternative is admitting you do not
know. He had the matching contempt for its mirror image, the will to believe, and he thought the
second had ruined more careers than the first.

\textbf{Four: the data must be collected by instruments that were already running.} Hynek's
late-life ambition was not a better witness form. It was a network of unattended monitoring
stations, recording continuously, generating a baseline of ordinary nights against which an
unusual one could be judged. He did not live to see it built. Section~\ref{sec:equipmenttools}
is, essentially, a report on how far his wish list has got.

\begin{funfact}
Hynek coined the term \emph{invisible college} for the loose, mostly anonymous network of
working scientists who were privately interested in UFO reports and professionally unable to
say so. He put the number in the low hundreds and he corresponded with a great many of them. The
single biggest change between his era and ours is that this college has become slightly visible
--- which is progress, and which also means the excuse of career risk has weakened considerably.
\end{funfact}

He also got things wrong, and it is worth saying which things, because a method is only credible
if you apply it to its author. The late Hynek entertained interdimensional and
quasi-metaphysical hypotheses on evidence that the 1955 Hynek would have binned without comment.
His close-encounter scale, discussed below, was extended by other hands into categories he never
sanctioned. And the organisation he founded to institutionalise his rigour outlived his ability
to enforce it. None of that touches the four rules. Judge the tools in this chapter by the rules,
not by the man's last decade.

% ---------------------------------------------------------------
\section{MUFON}\label{sec:mufon}

Start with the organisations, and start by noticing that all three of them are answers to
Hynek --- one built because he was not there, one built by him, and one built as if he had never
written anything.

The Mutual UFO Network was founded in 1969 --- as the Midwest UFO Network, in Illinois ---
in direct response to the closure of Project Blue Book. The logic was straightforward: if
the government has stopped collecting reports, somebody had better keep collecting them.
More than fifty years later MUFON operates the largest civilian UFO case database in the
world, with a public reporting form, a network of trained volunteer field investigators,
a monthly journal, and state and national directors.

I cite MUFON material in this book, including in Chapter~\ref{ch:language}, and I want to be clear about
why and with what reservations.

What MUFON does well is scale and continuity. Roughly 500 to 1,\!000 reports arrive per
month. The organisation runs a Field Investigator certification with an actual manual and
an actual examination, which is more procedural rigour than any other civilian body in the
field has managed to sustain. Cases are assigned dispositions, and --- this is the part
that matters --- a large majority are closed as identified or as insufficient data. That
denominator is precisely what Section~\ref{sec:doubtdrivenskepticism} says a serious field must publish.

What MUFON does badly is govern itself. It has been through repeated internal crises:
resignations of senior figures over the leadership's handling of allegations, controversy
over funding relationships with wealthy private donors, a 2013 episode in which a state
director's public statements caused a mass resignation of investigators, and persistent
friction between the organisation's scientific ambitions and its reliance on paid
conference appearances by people whose claims its own investigators would fail. The
journal's editorial standard varies wildly from issue to issue. Data quality is
self-reported and unaudited.

\begin{funfact}
MUFON's greatest scientific value may be entirely accidental. Because it has collected
reports continuously since 1969 using a broadly stable form, its database is one of the
few long-baseline datasets in existence --- which makes it useful for studying
\emph{reporting behaviour} even if every individual case were worthless. You can see the
Belgian and Phoenix triangle waves in it. You can see the arrival of camera phones. You
can see Chinese lanterns appear as a report category around 2005 and then fade. That is
real data about real human beings, whatever else is going on.
\end{funfact}

The practical advice is this. Use MUFON's database as an index, not as a source. A case
number tells you where to start looking; it does not tell you that anybody competent has
looked already. And apply the same standard to the two other bodies you will encounter:
NUFORC, the National UFO Reporting Center, which is essentially a raw unfiltered feed and
should be treated as one, and the now-defunct CUFOS founded by Hynek, whose case work was
the best the field has ever produced and whose archive is criminally underused.

The comparison is instructive and it is not flattering. Hynek's Center for UFO Studies, founded
in 1973 in Evanston, never had MUFON's volume and never wanted it. What it had was a scientific
board, case reports written to be checked, and a stated willingness to publish resolutions.
CUFOS investigated fewer cases than MUFON does in a quarter and produced more usable knowledge
than MUFON has in fifty years, for one reason: \hi{Hynek designed the organisation around the
residue rather than around the intake.} MUFON is optimised to receive reports. NUFORC is
optimised to publish them. Neither is optimised to close them, and closing them was the entire
point.

So when you use these databases, impose the missing layer yourself. Sort by observing conditions
before you sort by anything else --- Hynek's first rule --- and you will find the interesting
fraction of a thousand-case dump is usually under twenty. Those twenty are the only ones worth
your weekend.

% ---------------------------------------------------------------
\section{Classification of Common UFO Shapes}\label{sec:classificationcommonufos}

Before you can analyse reports you need a vocabulary that does not smuggle in conclusions.
The standard shape taxonomy is imperfect but it is what everybody uses, so learn it --- and then
learn why Hynek refused to build his system on it.

His own classification, published in \emph{The UFO Experience} in 1972, does not sort by shape at
all. It sorts by observing conditions, in six boxes: nocturnal lights, daylight discs,
radar-visual cases, and then the three close-encounter grades. Notice what that ordering does. A
nocturnal light is the weakest class of observation available, and it is also the most common ---
so the taxonomy tells you at a glance that the bulk of the corpus is intrinsically low-grade
evidence. A radar-visual case is rare and strong, because two independent detection channels
agreed. \hi{Hynek's categories rank your ability to know, not the object's appearance.} The shape
list below ranks appearance, which is why it is useful for indexing and dangerous for reasoning.
Keep both, and keep them in that order of authority.

\textbf{Discs.} The classic saucer. Dominant from 1947 until roughly 1980 and now
relatively rare in raw reports, which is itself informative --- if the phenomenon were
independent of culture you would not expect its shape to go out of fashion.

\textbf{Triangles.} Large, dark, silent, lights at the corners or vertices, slow. Dominant
from the mid-1980s onward. See Chapter~\ref{ch:historyarea51} on the F-117 and Tacit Blue, and note the timing.

\textbf{Spheres and orbs.} Now the most commonly reported category by a wide margin, and
the least useful, because a distant unstructured light is compatible with almost
everything: aircraft landing lights head-on, Venus, Starlink trains, drones, lanterns,
balloons, reflections.

\textbf{Cigars and cylinders.} Elongated, no wings, historically often reported as
``airships''. Prominent in the 1896--97 American airship wave, decades before powered
flight was public.

\textbf{Tic-tacs.} A recent addition, from the 2004 Nimitz encounter footage: a smooth
white oblong with no visible control surfaces, wings, or exhaust.

\textbf{Chevrons, boomerangs, and V-shapes.} The Phoenix Lights of 13 March 1997 are the
canonical case, and they are also two separate events conflated into one --- an early
formation of lights, plausibly aircraft, and a later stationary arc of lights over the
Estrella range which the Maryland Air National Guard identified as A-10 flare drops. The
second identification is solid. The first is not settled.

\figwrap[r]{0.36}{Reported morphologies tallied by shape. The categories are the reporter's, not
the object's --- which is exactly why the shape field alone constrains so little.}{https://commons.wikimedia.org/wiki/File:UFO_Sightings_by_Shape.svg}{fig:shapes}
Alongside shape you should record the observational parameters that actually constrain
anything: angular size (measured against a thumbnail at arm's length, not guessed in
feet), duration, bearing and elevation, apparent motion relative to fixed references,
sound, and whether the object was ever occluded by cloud or terrain. \hi{Distance and size
estimates from witnesses are worthless in isolation, because the human visual system
cannot judge either without a known reference.} A witness who says ``it was the size of a
football field'' is telling you about their impression, not about the object. A witness who
says ``it was about three finger-widths across and passed behind that water tower'' has
handed you geometry.

One caution about the whole taxonomy before you use it. A list of shapes invites you to assume
there is a shape --- a rigid object with one outline that a witness either read correctly or did
not. A stubborn minority of reports refuse that, describing outlines which change while under
observation, and the residual card in the taxonomy below exists because of them. This is the
Hynek residue in miniature: the cases that will not go in any box are the only ones that could
ever teach you anything, and the temptation to force them into a box is strongest precisely
there.
Section~\ref{sec:solidity} takes the problem seriously and asks what the assumption of solidity
actually rests on; the answer is one burned prospector at Falcon Lake
(Section~\ref{sec:falconlake}), and nothing else. Classify shapes, by all means. Do not let the
act of classifying persuade you that you know what you are classifying.

\begin{funfact}
Hynek's close-encounter scale, which everyone half-remembers from the Spielberg title, is
about \emph{proximity and effect}, not shape. CE1: observed within about 500 feet with no
physical interaction. CE2: physical traces --- scorching, vehicle interference, animal
reaction. CE3: occupants observed. CE4 (added later, not by Hynek) : abduction. CE5
(added later still, and not accepted by serious investigators): deliberate human-initiated
contact. The scale degrades in reliability at exactly the rate it gains in excitement.
\end{funfact}

% ===================== UFO SHAPE TAXONOMY =====================
% Seventy-two reported shapes in strict A--Z order, nine to a page over eight
% pages, on the same
% card system as the Alien and Cryptid rosters but three columns wide, and
% over the Earth plate rather than the starfield or the swamp. The last line
% of every card is the useful one: what the shape nearly always turns out to
% be once somebody does the work.

The shape catalogue below is the part of this chapter I expect to be argued
with, so let me say plainly what it is and what it is not. It is not a
zoology of craft. It is an index of \emph{reported outlines} --- what witnesses
drew, what investigators filed, and what each outline has historically
resolved into. Read it as a lookup table. When a report lands on your desk,
find the shape, read the last line of the card, and eliminate that first. If
the ordinary explanation survives the elimination, you have something worth
your evening.

\earthfieldon
\rosterwidethree

\begin{multicols}{3}[\rosterhead{Classification of Common UFO Shapes}]

\ufoshapecard{\gacorn}{Acorn}{Discoid; Kecksburg, 1965}{A rounded body with a
collar or banded base. Kecksburg is the case that carries the type, and it is
also the case that shows how fast a shape hardens: the drawings became more
consistent \emph{after} the story was retold, not before.}{A re-entering space
object, most plausibly a spent probe or booster section.}

\ufoshapecard{\gamoeba}{Amorphous}{Formless; the residual category}{No stable
outline: reported as shifting, folding, or changing between glances. Every
classification system needs a bin for the reports that will not sit still, and
this is mine. It is also the least useful bin.}{Cloud, smoke, atmospheric
refraction, or the limits of night vision.}

\ufoshapecard{\gbarrel}{Barrel}{Cylindrical; industrial-area sightings}{A short
banded cylinder, sometimes reported as ribbed or seamed. Reports cluster near
ports, refineries and rail yards --- that is, near the places barrels
are.}{Industrial plant, a storage tank, or a shipping container under a crane.}

\ufoshapecard{\gbell}{Bell}{Discoid; various, 1940s onward}{Flared at the base,
narrowing to a crown. It appears in wartime and post-war European reports and
later in the exotic-engine literature, which is where the trouble starts: the
shape is now inseparable from claims about propulsion that nobody has
demonstrated.}{An illuminated tethered balloon, or an aircraft nose-on at
night.}

\ufoshapecard{\gblacktri}{Black Triangle}{Angular; Phoenix Lights, 1997}{The
night variant: unlit hull inferred from the pattern of lamps against the sky.
Phoenix showed exactly how this class fails --- one event was flares, another
was aircraft, and the \emph{shape} was supplied by the witnesses.}{Flares on
parachutes, an aircraft formation, or occlusion of stars misread as a hull.}

\ufoshapecard{\gboomerang}{Boomerang}{Angular; Hudson Valley, 1983--84}{A
V-shaped outline with a shallow apex, reported as gliding without sound. The
Hudson Valley wave was eventually traced in large part to a group of pilots
flying in formation for their own amusement, and they said so.}{Ultralight
formation flying, or migrating birds lit from below by city glow.}

\ufoshapecard{\gfractal}{Branching Form}{Formless; very rare}{A repeating
branched structure, reported chiefly in altered-state and abduction accounts
rather than in field sightings. The reporting context is the finding.}{Entoptic
phenomena --- form constants generated by the visual system itself.}

\ufoshapecard{\gbutterfly}{Butterfly}{Organic; sparse}{Paired lobed wings,
sometimes reported with markings. It is the class most likely to be an insect
near the lens, and the least likely to be anything at altitude.}{An insect
close to the camera, or a bird photographed mid-wingbeat.}

\ufoshapecard{\gcapsule}{Capsule}{Cylindrical; re-entry and test
windows}{Rounded at both ends, reported as descending steadily and sometimes as
under a canopy. Where a canopy is mentioned the case usually closes
immediately.}{A returning crew or cargo capsule, or a high-altitude test
article under parachute.}

\ufoshapecard{\gchain}{Chain of Lights}{Formation; Starlink era, post-2019}{A
regular line of evenly spaced lights holding station relative to each other.
This class exploded after 2019 and the cause is now so well established that
it has become a useful test of a reporting system's competence.}{A satellite
constellation train shortly after launch.}

\ufoshapecard{\gcigar}{Cigar}{Cylindrical; the 1896--97 airship flap}{An
elongated body, blunt at both ends, reported as huge and slow. It dominates
the pre-aviation record because it was the plausible future craft of the day.
Shape reports track the technology the public expects.}{An airship or
dirigible, a contrail lit end-on, or an aircraft seen along its axis.}

\ufoshapecard{\gdisc}{Classic Disc}{Discoid; Kenneth Arnold, 1947}{The
foundation shape of the whole subject, and an accident of language: Arnold
described the \emph{motion} as a saucer skipping on water, and the press gave
the world a saucer-shaped object. Every drawing filed afterwards was already
contaminated by that headline.}{A lenticular cloud, a weather balloon at
altitude, or the witness reproducing a picture they have already seen.}

\ufoshapecard{\gcloudcigar}{Cloud Cigar}{Composite; French wave, 1954}{An
elongated form with a cloudy, ill-defined envelope, sometimes reported
emitting smaller objects. Aimé Michel built a whole theory on these, and the
theory has aged considerably worse than the observations.}{A contrail or
lenticular formation, with aircraft nearby read as ejected craft.}

\ufoshapecard{\gcluster}{Cluster}{Formation; wave reports}{Several bodies
moving in loose company, described collectively as one phenomenon. Clusters
are where witness counts inflate fastest, and where the number reported is
least likely to be the number present.}{Balloon releases, sky lanterns at a
festival, or a drone display.}

\ufoshapecard{\gcomet}{Comet Form}{Luminous; spiral and plume events}{A bright
nucleus with a broad persistent tail, sometimes reported as spiralling. The
Norwegian spiral of 2009 was the making of this class and it was a failed
missile stage, confirmed and admitted.}{A rocket-stage propellant dump, or a
spiralling failed launch seen at range.}

\ufoshapecard{\gcone}{Cone}{Angular; nose-on and re-entry reports}{A circular
base drawn to an apex, reported both in flight and on the ground. It is what a
great many other shapes look like when they are pointed at you, so treat cone
reports as aspect-angle problems.}{A re-entry capsule or nose fairing, or an
aircraft seen nose-on.}

\ufoshapecard{\gcoolie}{Coolie Hat}{Discoid; Papua and Java reports, 1950s}{A
shallow cone over a flat base, reported chiefly from South-East Asia and the
Pacific in the early years. The regional concentration tells you it is at
least partly cultural: the witness reaches for a familiar silhouette.}{Cloud
form, or an aircraft seen from below with wing and fuselage merged.}

\ufoshapecard{\gcrescent}{Crescent}{Curved; wartime and Nazi-craft claims}{A
swept arc with concave trailing edge. It attracts the worst literature in the
subject, and the shape's own history --- a real German programme of flying
wings --- is more interesting than anything written about it since.}{A flying-wing
aircraft, or the moon's crescent through broken cloud.}

\ufoshapecard{\gcross}{Cross}{Angular; European reports, 1960s--70s}{Two
intersecting bars of light, reported as maintaining orientation. The reported
rigidity is the diagnostic feature and also the least reliable, because
witnesses are poor judges of rotation at distance.}{An aircraft head-on with
landing lights and wingtip strobes.}

\ufoshapecard{\gcube}{Cube in Sphere}{Geometric; contemporary, post-2015}{A
cubic form apparently inside a transparent shell, a description that arrived
late and spread fast online. Note the date. New shapes still enter the
catalogue, and they enter through the internet, not the sky.}{Composite or
edited imagery; the class has essentially no pre-digital record.}

\ufoshapecard{\gvoid}{Dark Mass}{Formless; nocturnal reports}{An area of sky
blacker than the sky, defined only by the stars it blocks. This is inference,
not observation, and it is the single easiest class to generate by accident
with a passing cloud.}{Cloud occluding starfield, or a large unlit balloon.}

\ufoshapecard{\gdelta}{Delta}{Angular; military airspace reports}{A swept
triangular planform with a notched trailing edge --- that is to say, a delta
wing. When a shape class matches a production aircraft this exactly, the
burden of proof sits with whoever claims it is not one.}{A delta-winged
aircraft; the type has been in service since the 1950s.}

\ufoshapecard{\gdiamond}{Diamond}{Geometric; Rendlesham Forest, 1980}{Two
opposed pyramids, reported as a lit lozenge low among trees. Rendlesham is the
case I would most like to be real and the case where the lighthouse timing
matches the testimony best.}{A lighthouse beam through trees, a bright star
low on the horizon, or re-entry debris.}

\ufoshapecard{\gdome}{Dome}{Discoid; ground-contact reports}{A hemisphere,
usually reported on or near the ground rather than in flight. It is the
landed-object shape, and it has the highest rate of physical-trace claims
attached to it, which at least makes it testable.}{Farm equipment, a grain
silo top, a tent, or a radome seen at distance.}

\ufoshapecard{\gdomedisc}{Domed Disc}{Discoid; Rex Heflin, 1965}{A disc with a
raised cupola, usually reported as glazed or windowed. It is the most drawn
variant in the archives because it is the most cinematic, and because a dome
gives the witness somewhere to put the occupants they have not seen.}{Hubcaps,
lens flare on a bright object, and outright models suspended on thread.}

\ufoshapecard{\gdumbbell}{Dumbbell}{Paired; scattered, all decades}{Two masses
joined by a slimmer waist. Whenever a report describes two connected objects,
ask first whether the connection was seen or inferred; it is nearly always
inferred.}{Two objects at different distances aligned by chance, or a tethered
pair of balloons.}

\ufoshapecard{\gegg}{Egg}{Ovoid; South American reports, 1950s--70s}{An ovoid
with one broader end, often reported as landed with occupants nearby. The
occupant claims come attached to the shape rather than the reverse, which
suggests the shape is doing narrative work.}{A hot-air balloon envelope
partially inflated, or a tank or vessel seen in a field.}

\ufoshapecard{\gfincyl}{Finned Cylinder}{Cylindrical; Cold War intrusion
reports}{A tube with tail surfaces, reported chiefly near military ranges and
coastlines. It is the shape most likely to have a boring answer, because the
places it is reported from are the places such objects are actually
launched.}{A test missile, sounding rocket, or target drone.}

\ufoshapecard{\gfireball}{Fireball}{Luminous; ubiquitous}{A brilliant
moving light with a visible head and short tail, usually reported for seconds
rather than minutes. It is the highest-volume report class worldwide and the
one with the highest resolution rate, which should tell you
something.}{A bolide meteor, or re-entering satellite debris burning up.}

\ufoshapecard{\ggyro}{Gyroscopic}{Composite; sparse}{Nested rotating rings
around a central body, usually reported by a single witness at length. Detail
of this order from one observer and none from anybody else is a reliable
warning sign.}{Confabulation, or a defocused light with rotational camera
movement.}

\ufoshapecard{\gnut}{Hexagonal Nut}{Geometric; sparse, modern}{A six-sided
plate with a central aperture, reported as flat and edge-lit. It appears
almost only in still photographs, never in multi-witness sightings, and that
asymmetry tells you where to look for the cause.}{A lens or sensor artefact
shaped by a hexagonal aperture blade set.}

\ufoshapecard{\ghourglass}{Hourglass}{Paired; sparse, modern}{Two cones meeting
at a narrow waist, sometimes reported as pulsing in step. Rare enough that
every entry in the class can be examined individually, and most of them are
photographs rather than sightings.}{An optical artefact --- an out-of-focus
light source with an iris or coma aberration.}

\ufoshapecard{\gicecream}{Ice Cream Cone}{Composite; ground-trace cases}{A
sphere seated in a cone, most often reported as landed on legs or a tripod.
The landing-gear detail is doing a lot of work here, and it is the detail
witnesses most often add after the fact.}{Farm or survey equipment, a
floodlight on a stand, or a hoax construction.}

\ufoshapecard{\gbowl}{Inverted Bowl}{Discoid; ground-level reports}{A shallow
concave form, reported open-side-down and usually stationary. It is the
second of the two landed-object shapes and it shares the first's habit of
turning out to be agricultural.}{A radome, a satellite dish, or a
polytunnel end seen at dusk.}

\ufoshapecard{\gjelly}{Jellyfish}{Organic; Iraq thermal footage, 2018}{A domed
body with trailing filaments. The declassified thermal clip revived this
class, and the honest answer is that a thermal sensor renders a mundane object
in an unfamiliar way and the eye supplies tentacles.}{Thermal blooming around
a balloon or drone, with sensor smear read as trailing limbs.}

\ufoshapecard{\gkite}{Kite Form}{Angular; daylight reports}{A four-sided
symmetrical outline, tapering below, reported as hovering and yawing. There is
a reason the shape shares its name with the object, and the object is
common.}{A kite, or a bird soaring with wings and tail spread.}

\ufoshapecard{\glampshade}{Lampshade}{Composite; scattered}{A truncated cone,
narrow above and flared below, often reported as glowing from underneath. The
underside glow is the detail that most often turns out to be the light source
itself.}{A streetlamp or floodlight seen through mist, or an illuminated
balloon.}

\ufoshapecard{\glens}{Lenticular}{Discoid; McMinnville, 1950}{Symmetrical,
thicker at the centre, tapering to an edge --- which is precisely what a cap
cloud looks like when it forms over a ridge. The name of the meteorological
phenomenon and the name of the shape class are the same word for a
reason.}{Altocumulus lenticularis, sun through thin cloud, or a lens
artefact.}

\ufoshapecard{\gorb}{Luminous Orb}{Spherical; ubiquitous, all decades}{The
single commonest report in every national archive. A point of light with
apparent volume, no structure, and no scale. It is also the class with the
highest proportion of cases closed by simply asking what was in the sky that
night.}{Venus, Jupiter, a sky lantern, a drone, or a distant landing light.}

\ufoshapecard{\gmanta}{Manta}{Organic; modern reports}{A broad winged form with
a swept trailing edge, reported as undulating rather than banking. Undulation
is the reported feature that most reliably indicates a living thing.}{A large
bird gliding, or a flying-wing airframe at distance.}

\ufoshapecard{\gsphere}{Metallic Sphere}{Spherical; Coyne helicopter case,
1973}{A featureless reflective ball, frequently described as \emph{seamless}
and as showing no exhaust. The absence of features is the whole problem: a
sphere gives you nothing to measure, no aspect angle, and therefore no
distance.}{A high-altitude research or weather balloon, which is spherical,
reflective and silent by design.}

\ufoshapecard{\gslab}{Monolith}{Geometric; Utah and copycats, 2020}{A tall
rectangular slab, usually reported standing rather than flying. The 2020
installations were art, they were admitted as art, and they still generated
sincere anomaly reports for weeks afterwards.}{Sculpture, a survey marker, or
a communications cabinet in poor light.}

\ufoshapecard{\gmushroom}{Mushroom}{Composite; scattered, post-1945}{A domed
cap over a narrower stem. The date matters: the shape enters the record
immediately after the mushroom cloud enters public consciousness, which is a
pattern you will see again and again in this list.}{Cloud formation, a water
tower, or a lit structure with the base in shadow.}

\ufoshapecard{\goct}{Octahedron}{Geometric; scattered modern reports}{Faceted,
angular, described as tumbling or rotating with a repeating glint. The
repeating glint is the giveaway and also the only measurable thing in the
report: time the period.}{A tumbling satellite or rocket body catching the
sun.}

\ufoshapecard{\gpillar}{Pillar of Light}{Luminous; cold-climate reports}{A
vertical column of light with no visible source at either end, reported
standing still for minutes. It is spectacular, it photographs well, and it has
a completely settled meteorological explanation.}{A light pillar --- ground
lighting reflected from ice crystals in still, cold air.}

\ufoshapecard{\gplasma}{Plasma Ball}{Spherical; Hessdalen, from 1981}{A
glowing sphere with a visible corona or flicker, reported in geologically
active valleys. Hessdalen is the one place where instruments were taken to the
phenomenon rather than the other way round, which is why it is worth
respecting.}{Ionised gas from piezoelectric or tectonic stress; ball lightning
remains poorly characterised.}

\ufoshapecard{\gpoint}{Point Source}{Luminous; the base case}{No shape at all:
a light with no resolvable extent, no structure, and no way to establish
distance. It is the honest terminus of any shape taxonomy and, by volume, the
largest single class in every archive there is.}{Aircraft, satellites,
planets, or drones --- an unresolved light, and nothing more.}

\ufoshapecard{\gpyramid}{Pyramid}{Geometric; Russell Island and naval
footage}{A triangular-faced solid with a defined apex, more often reported
stationary than moving. It photographs badly and reads triangular from almost
every angle, which flatters the drawing.}{Defocused point lights --- a
triangular bokeh from a camera iris will do it.}

\ufoshapecard{\gring}{Ring}{Annular; \emph{ovnis en anneau}, France 1970s}{A
thin circle or halo with an apparently empty centre. French investigators
filed a good number of these, and the honest reading is that a ring is what
your eye makes of a defocused bright point.}{A camera artefact, a lit
helicopter rotor arc, or a dispersing contrail.}

\ufoshapecard{\grocket}{Rocket Form}{Cylindrical; twilight launch
windows}{Pointed body with fins and, crucially, a visible plume. This is the
one class that resolves almost perfectly: check the launch schedules for the
hour and the case is normally closed before dinner.}{An actual rocket launch
seen at range, with the plume illuminated by high sunlight.}

\ufoshapecard{\grod}{Rod}{Elongated; camera-artefact era, 1990s}{A long
segmented streak, found in photographs and video and almost never seen by eye.
José Escamilla built a career on these and the answer was available from the
beginning to anyone who understood shutter speed.}{Motion blur on a nearby
insect --- a rolling-shutter and exposure artefact.}

\ufoshapecard{\gsaturn}{Saturn Form}{Discoid; Trindade Island, 1958}{A central
body with an encircling rim or ring. The Trindade photographs made the type
famous and have since been picked apart by everyone who has handled the
negatives. Note that the shape is easy to build and easy to throw.}{A thrown
or suspended model, or a bright object smeared by camera shake.}

\ufoshapecard{\ginsect}{Segmented Body}{Organic; close-range reports}{A jointed
body with limbs or projections, reported only at close range and in poor
light. Close range should mean better data; here it reliably means worse,
because it means one startled witness.}{An insect or bird near the observer,
with scale badly misjudged.}

\ufoshapecard{\gsheet}{Sheet}{Formless; sparse}{A broad flat expanse, reported
as rippling or flexing while it moves. Rigid objects do not ripple, so the
report is either of something flexible or of something being misjudged.}{A
sheet of plastic or tarpaulin lifted on a thermal, or aurora at low
latitude.}

\ufoshapecard{\gsombrero}{Sombrero}{Discoid; Latin American reports}{A wide
flat brim with a raised crown --- the disc family's regional variant, and named
for a hat again. Note how often the catalogue names a shape after whatever the
witness had to hand.}{A lenticular cloud, or a disc-shaped object photographed
from near its plane.}

\ufoshapecard{\gspindle}{Spindle}{Elongated; high-altitude reports}{Tapered to
a point at both ends and reported at great height and speed. The height claim
is the weak link: without a second observer or a radar track, altitude is a
guess dressed as data.}{A high-altitude balloon seen end-on, or a contrail
lit at both ends.}

\ufoshapecard{\gtop}{Spinning Top}{Composite; Brazilian and Japanese
reports}{A conical body with a spindle above, described as rotating on its
axis. It is a toy shape, and I mean that literally: several photographic cases
in this class were resolved to actual toys.}{A thrown or suspended model, or a
lit spinning object close to the camera.}

\ufoshapecard{\gwheel}{Spoked Wheel}{Annular; \emph{Ashtar} wheel, Gulf of
Oman}{A rotating disc with radial spokes of light, reported almost exclusively
at sea. It is one of my favourite entries because it has a beautiful
explanation that most people find less interesting than the mystery.}{Marine
bioluminescence in an internal wave field, seen from the deck.}

\ufoshapecard{\gsquare}{Square}{Geometric; sparse, all decades}{A flat
quadrilateral, reported edge-on as a line and face-on as a panel. It is rare,
and rarity here is a clue: genuinely novel geometry should not be rare if it
is flying about.}{A kite, a lifted tarpaulin or advertising panel, or a
solar-panel glint.}

\ufoshapecard{\gstar}{Star Form}{Luminous; contemporary reports}{A radiating
point with distinct arms, reported as brilliant and often as colour-changing.
This is the class where astronomy closes the file fastest: an app, a star chart
and a timestamp will normally do it.}{Sirius or another bright star scintillating
low in the atmosphere.}

\ufoshapecard{\gteardrop}{Teardrop}{Ovoid; sparse, all decades}{Rounded at one
end, drawn to a point at the other, usually reported as descending. It is the
shape a falling luminous object appears to take because the eye adds a tail to
motion it cannot resolve.}{A meteor or bolide, or burning debris on a ballistic
path.}

\ufoshapecard{\gtee}{Tee}{Angular; sparse}{A crossbar over a stem, typically
reported near airfields and installations. When a shape is reported almost
exclusively where aircraft manoeuvre, treat the location as part of the
evidence.}{An aircraft in a banking turn, or a crane or mast against the sky.}

\ufoshapecard{\gtictac}{Tic Tac}{Cylindrical; USS Nimitz FLIR, 2004}{A smooth
rounded oblong with no wings, no rotor and no plume. This is the modern
flagship case, and my position on it is unchanged: the video is genuine, the
interpretation is not settled, and the sensor is not a witness.}{A distant
aircraft with parallax and glare effects on an infrared sensor.}

\ufoshapecard{\gtophat}{Top Hat}{Composite; European reports, 1954 wave}{A
cylinder set on a flat brim. The 1954 French wave produced a great many of
these, along with the best-documented mass panic in the field's history and
very little in the way of physical evidence.}{An illuminated water tower or
silo, or a lantern seen against low cloud.}

\ufoshapecard{\gtoroid}{Toroid}{Annular; smoke-ring and vortex reports}{A ring
with visible thickness, sometimes reported as rotating about its own axis. It
is a shape that fluid dynamics produces cheaply, which is exactly why it needs
no exotic explanation.}{A vortex ring from an explosion or industrial vent, or
a volcanic gas ring.}

\ufoshapecard{\gtorpedo}{Torpedo}{Cylindrical; naval and coastal sightings}{A
rounded nose tapering to a blunt tail, usually reported low over water and
often as transmedium. The transmedium claim is the interesting part and the
part never captured by more than one sensor at a time.}{A submarine masthead,
a whale surfacing, or an aircraft seen from an unusual angle.}

\ufoshapecard{\gclear}{Transparent Craft}{Formless; rare, modern}{Reported as a
visible distortion or refraction rather than a solid --- ``like heat haze with
edges''. Where a claim is unfalsifiable by construction, it is not evidence,
whatever else it is.}{Heat shimmer, a window or windscreen reflection, or a
soap-film artefact on the lens.}

\ufoshapecard{\gtrapezoid}{Trapezoid}{Geometric; sparse}{A flat form with
unequal parallel edges, reported as tilting. In practice it is nearly always a
rectangle seen in perspective, which is a thing witnesses are demonstrably bad
at correcting for.}{A kite, banner, or panel; or a rectangular object seen
obliquely.}

\ufoshapecard{\gtriangle}{Triangle}{Angular; Belgian wave, 1989--90}{A
flat-bottomed triangle with a light at each vertex, silent, slow, very large.
The Belgian wave produced hundreds of consistent reports and no photograph
anybody can work with, which is the recurring pattern of this field.}{Formation
flying, a stealth or experimental airframe, or three lights read as one
object.}

\ufoshapecard{\gvform}{V Formation}{Formation; Lubbock Lights, 1951}{Lights in
a V, reported as a rigid object rather than a group. Lubbock was resolved to
birds reflecting new street lighting, and the resolution was published, and
almost nobody remembers it.}{Migrating birds catching artificial light from
below.}

\ufoshapecard{\gwedge}{Wedge}{Angular; testing-corridor sightings}{Asymmetric,
thick at one end and knife-edged at the other. Sightings cluster tightly
around known flight-test ranges, and the reports get more detailed the closer
the witness is to a fence they cannot see over.}{A prototype or classified
airframe under test, which is not the same claim as a spacecraft.}

\ufoshapecard{\gwye}{Wye}{Angular; South American reports}{Three arms from a
common hub, described as slowly rotating. It is the shape you would expect if
three lights in loose formation were read as one rigid object, which is
normally what has happened.}{Three lights in formation, or a rotating
navigation beacon.}

\end{multicols}

\rosternarrow
\earthfieldoff


% ---------------------------------------------------------------
\section{Photographic Analysis}\label{sec:photographicanalysis}

I told you in Section~\ref{sec:ai} that photographic evidence is finished as a category. It
is, for anything captured from now on. But the historical corpus still needs analysing, and
the techniques are worth knowing because they generalise to any image claim.

Hynek's position on photographs was harder than most sceptics'. He held that a photograph, in
isolation, is close to worthless \emph{as evidence about an object}, because it discards exactly
the information you need: it cannot fix distance, it cannot fix size, and it cannot establish that
anything was there beyond the film. What it can do is corroborate a witness whose account was
secured independently. That is a supporting role, not a starring one, and the field has had it
backwards for seventy years --- building cases on images and treating the humans as colour.

Start with the metadata, and start there because it is where most hoaxes die. EXIF data
gives you camera model, timestamp, focal length, exposure, and often GPS. Cross-check the
claimed time and place against sun and moon position, weather records, and aircraft
tracking archives. A great many cases resolve in ten minutes at this stage, and the most
common resolution is that the timestamp is wrong and the object is the Moon or Venus.

Then the physics of the image itself. Four questions do most of the work.

\emph{Is the object in focus, and does its focus match the background?} A nearby small
object photographed with a background-focused lens will show characteristic blur. An
object sharper than the landscape behind it was not at the distance claimed.

\figwrap[r]{0.38}{One of the two Trent photographs, McMinnville, Oregon, 11 May 1950. Note the overhead wires at the top of the uncropped frame.}{https://commons.wikimedia.org/wiki/File:McMinnvilleUFO.jpg}{fig:trent}
\emph{Does the lighting match?} This is the single most reliable discriminator for both
physical models and digital composites. The direction, colour temperature, and hardness of
light on the object must match the scene. Shadow direction on the object must be consistent
with shadows on the ground.

\emph{Is there motion blur, and is it consistent?} An object claimed to be moving at high
speed in a 1/60-second exposure should be smeared. A crisp object in a frame where the tree
line shows camera shake was moving with the camera --- which means it was attached to it, or
very close.

\emph{What is the angular size and what does it imply?} With focal length, sensor size, and
pixel measurement you can compute the object's angular size precisely. Angular size plus
any distance assumption gives physical size. This is how the Trent/McMinnville photographs of 1950 (see Figure~\ref{fig:trent}) --- still among the best-regarded in the field --- have been argued over for
seventy years: the debate turns entirely on whether brightness measurements of the
underside indicate a distant large object or a nearby small one hung from the overhead
wires that were visible in the original negatives and cropped out of most reproductions.

\begin{funfact}
The most famous UFO photographs almost all have a cropping problem. The McMinnville images
originally showed power lines. The 1967 Alex Birch photograph was admitted a hoax by its
teenage author. The 1990 Petit-Rechain Belgian triangle photograph --- reproduced in
hundreds of books as the best triangle image in existence --- was confessed as a fake in
2011 by the man who made it, using polystyrene and a lamp. In each instance the field had
spent decades analysing the pixels of an image whose provenance nobody had audited. Chain
of custody beats image processing every time.
\end{funfact}

For video, add frame-by-frame examination for compression artefacts, check whether apparent
object motion correlates with camera motion, and be extremely careful with footage from
infrared and gimballed targeting systems. The 2004 and 2015 US Navy videos are genuine
recordings of genuine sensor events, and much of what looks anomalous in them --- the
sudden ``rotation'' of the Gimbal object, the apparent zip of the GoFast object across the
water --- has plausible explanations in the behaviour of the gimbal mechanism and in
parallax. ``Plausible'' is not ``established''. Neither is ``anomalous''. The honest
reading of those videos is that the sensor recorded something the operators could not
identify, and that the recording alone cannot resolve size, distance, or speed. Everything
beyond that is inference, including the exciting parts.

If you want Hynek's verdict on the whole category, it is in his handling of the Trent
photographs. He rated them among the best in existence, said so in print --- and still would not
build an argument on them, because the couple's account and the negatives together could not
exclude a small object near the camera. A man who tells you his best photograph is not good
enough is a man whose other judgements are worth reading.

% ---------------------------------------------------------------
\section{Investigative Techniques}\label{sec:investigativetechniques}

Before the procedure, a word about what an investigator is actually collecting, because the
list is longer and stranger than outsiders assume. This evidence includes but is not limited
to physical debris, radiation exposure, records from radar stations, photographs and drawings,
eye witness testimony, video and audio, unexplainable death, missing persons cases, damage to
surrounding area, traces of mysterious toxins, implantations, and yes EVEN DISAPPEARING BABIES
FROM THE WOMB OF THEIR MOTHER!!! Every category on that list has a chapter or a section behind
it later in this book, and every category on that list is handled the same way here.

Here is the procedure, in order, and the order is the whole point. Do not skip forward
because a case looks interesting.

Every step below is Hynek's, whether the people who use it know that or not. He built the
sequence in the field between 1948 and 1972, out of the cases he had personally mishandled, and
the reason it looks obvious now is that he did it first.

\textbf{1. Secure the account before you contaminate it.} Get the witness to write or
record their narrative before you ask a single specific question, and before they have read
anything about similar cases. Every question you ask plants something. ``Did you see any
windows?'' produces windows. This is not a hypothetical risk; it is the documented
mechanism behind most of the abduction literature, which we will get to in
Chapter~\ref{ch:abduction}. Use open prompts only: ``What happened next?''

\textbf{2. Reconstruct the sky.} Before forming any opinion, establish what was actually
up there. Astronomical software for planetary and stellar positions, ISS and satellite pass
predictions, Starlink train schedules, aircraft tracking archives, airport and military
notices to airmen, weather records including cloud base and inversion layers, and the
location of any nearby launch or test range. Do this for the exact stated time, and then do
it again for an hour either side, because witnesses are unreliable about time.

\textbf{3. Interview separately.} Multiple witnesses must be interviewed apart and before
they confer. Convergence between independently obtained accounts is meaningful. \hi{Convergence
after they have talked to each other is nothing at all.}

\textbf{4. Go to the site.} Photograph the witness's sightline from their stated position
with a known focal length so angular sizes can be recovered later. Measure bearings with a
compass, correcting for magnetic declination. Note every light source, tower, reflective
surface, and flight path. An astonishing number of cases resolve when you stand where the
witness stood and look where they looked.

\textbf{5. Collect physical evidence properly or not at all.} Soil samples from claimed
landing sites need controls --- same depth, same soil type, from outside the affected area,
collected at the same time. \hi{Without controls a laboratory result means nothing whatsoever.}
Metal fragments go to an independent laboratory with no knowledge of the case's origin, as
Section~\ref{sec:deconstructingpeerreview} argued. Photograph everything in place, with scale, before touching it.

\textbf{6. Write the negative result.} If the case resolves, publish it as resolved, with
the same care you would have given a positive. This is the step everybody skips and it is
the step that makes the difference between a field and a hobby.

Hynek is the proof of the last step in both directions. He spent twenty years as the man who
wrote the negative results --- so effectively that his own name became a joke for it --- and the
reason his residue argument landed in 1972 was that nobody could accuse him of having failed to
try. \hi{The credibility to say ``this one does not resolve'' is earned entirely by the hundreds
of times you said ``this one does.''} An investigator with no debunkings on their record has no
standing when they announce a mystery, and the field is full of such people.

One technique deserves an explicit warning: do not use hypnosis, and do not accept evidence
derived from it. I will make the full case in Section~\ref{sec:hypnosis}, but the summary
is that hypnosis does not recover memory, it manufactures confident memory, and once a
witness has been hypnotised their original account is permanently compromised. An
investigator who hypnotises a witness has destroyed their own evidence.

Finally, a note on temperament. Most of the people who report these things are decent,
sober, and frequently reluctant --- many are afraid of being mocked, and quite a few have
something professional to lose. Treating a witness with respect is not a courtesy, it is
methodology: a witness who feels judged will either shut down or start performing, and
neither gives you usable data. You can tell someone their sighting was Venus without
telling them they are stupid. I have had to do it more times than I can count, and the ones
who thank you for it are the reason this is worth doing.

That temperament is also Hynek's, and it is the part of him worth copying most. He is the only
figure in this history who publicly changed his mind in the direction that cost him something,
and he did it without ever deciding the witnesses had been fools. \hi{He separated the quality of
a person from the quality of their evidence} --- a distinction the debunkers and the believers both
collapse, in opposite directions, and the single most useful habit in this chapter.

% ---------------------------------------------------------------
\section{Equipment \& Tools}\label{sec:equipmenttools}
\label{sec:equipment}

Instrumentation is where ufology most wants to be taken seriously, and it is where the field
most often embarrasses itself. Most of the
kit described in this section is ordinary professional equipment, well characterised, sold to
geologists and firefighters and radiation safety officers who use it competently every day.

\actualq[problem]{that the instruments are bad.}{that a reading means nothing without a
baseline --- and the field has an extraordinary appetite for readings without baselines.}

So before any of the hardware, the rule that governs all of it. \hi{An instrument does not
detect anomalies. It produces numbers. Whether a number is anomalous is a claim about the
distribution of numbers you would have got anyway} --- at that site, with that device, in that
weather, in the absence of anything unusual. If you did not measure that distribution, you have
not detected anything. You have taken a reading.

This is not pedantry, and it is not scepticism for its own sake. It is the difference between
the Hessdalen programme, which has produced publishable data for four decades, and the vast
quantity of field video in which somebody waves a meter at a patch of grass and announces a
hotspot. Same instruments. Opposite epistemic value.

\Needspace*{0.40\textheight}
\subsection{Geiger Counter}

\figwrap[r]{0.34}{A handheld Geiger--M\"uller survey meter of the type carried on field investigations. The reading it gives you is meaningless without a matched background measurement taken away from the claimed site.}{}{fig:geiger}
A Geiger--M\"uller counter detects ionising radiation --- alpha, beta and gamma, depending on
the tube and window --- by counting discharge events in a gas-filled tube. It is cheap, robust,
and the single most misused instrument in the field.

What it is good for is establishing that a claimed landing site is \emph{not} radioactive. That
is a real and useful result, and it is the result you will get essentially every time. Natural
background varies by a factor of several depending on altitude, geology, building materials and
recent rainfall, and a granite outcrop or a load of monazite sand in the local aggregate will
read higher than the surrounding field for entirely boring reasons.

The canonical illustration is Rendlesham Forest, which I take up in
Section~\ref{sec:rendlesham}. Lieutenant Colonel Halt's memo records beta/gamma readings of
roughly 0.1 milliroentgen in the ground depressions and 0.07 mR/hr peaks, against a stated
background of 0.03 to 0.04. This has been quoted for forty years as elevated radiation at a
landing site. Health physicists who have looked at it are close to unanimous: with that
instrument, an AN/PDR-27, at that scale, those numbers are unremarkable variation and well
inside the range you would expect from soil, moisture and instrument response.
\hisoft{The reading was real. The significance was invented afterwards.}

\begin{funfact}
\funfactlead{The wrong instrument, correctly used.} The AN/PDR-27 was built to find battlefield
contamination at levels that would eventually kill you, not to do trace environmental survey
work. Asking it to resolve a 0.03 mR/hr difference is like weighing a letter on a lorry
weighbridge. Halt was not being dishonest; he was using the meter that was in the cupboard,
which is what everybody does.
\end{funfact}

Protocol, if you use one: log continuously rather than spot-reading, take at least thirty
minutes of background at matched geology outside the claimed area, record the tube type and
window material, note the weather including recent rain and radon conditions, and publish the
raw count series rather than a maximum. A maximum with no distribution behind it is not a
measurement.

\Needspace*{0.40\textheight}
\subsection{Ground-Penetrating Radar}

\figwrap[l]{0.34}{A ground-penetrating radar cart in use. The instrument maps dielectric discontinuities beneath the surface; it does not tell you what made them, or when.}{}{fig:gpr}
GPR pushes radio pulses into the ground and times the reflections, mapping dielectric
discontinuities --- buried voids, disturbed soil, backfilled trenches, rock. It is standard
archaeological and forensic equipment and it works, within limits set by soil conductivity: dry
sand is excellent, wet clay is nearly opaque.

Its one serious outing in this field was the 2002 excavation of the claimed Roswell debris
field on former Foster Ranch land, funded by the Sci-Fi Channel and carried out by a University
of New Mexico team under William Doleman. The remote sensing did find an anomaly: a shallow
linear feature consistent with something having gouged the ground. Then the team did the thing
that distinguishes archaeology from television. They excavated it and ground-truthed it. The
feature could not be shown to be anything other than natural erosional and geological
patterning, and no debris attributable to 1947 was recovered.

That is a model result, and it was widely reported as a failure, which tells you more about the
audience than the method. \hi{GPR cannot tell you what made a feature, only that a feature is
there.} Anything that ever disturbed the soil --- a plough furrow, a cattle track, a fence line,
a drainage channel, the 1947 clean-up itself if it happened --- produces the same class of
signature. The instrument narrows down where to dig. The spade decides.

\Needspace*{0.40\textheight}
\subsection{LIDAR}

\figwrap[r]{0.36}{A terrestrial laser scanner surveying a bridge. Used this way --- fixed station, known geometry --- LIDAR is sound. Pointed at empty air, a return gives you range and intensity, not identity.}{}{fig:lidarscan}
LIDAR ranges on laser returns to build a three-dimensional point cloud, and it has two distinct
uses here. The first is terrain and structure survey, which is what it is doing in
archaeological work and at Skinwalker Ranch: mapping ground surfaces, finding subsurface
disturbance signatures under vegetation, documenting a site geometry precisely enough that a
witness sightline can be reconstructed months later. That application is sound and I recommend
it.

The second use is atmospheric, and a LIDAR return from apparently empty air tells you there is
something there with a scattering cross-section --- aerosol, dust, ice, insects, a bird,
precipitation. This is where the claims get made, and it is where the physics needs stating
plainly. A return gives you range and intensity, not identity. Sequential returns give you a
track, and a track from a single station gives you angular motion, not velocity, unless you
already know the range independently at every point. \hisoft{Two units on a known baseline are
worth ten times one, and the field almost never deploys two.}

\Needspace*{0.40\textheight}
\subsection{Spectrometer}

\figwrap[l]{0.36}{A portable X-ray fluorescence spectrometer, one of the family of instruments that answers the question everything else dodges: what is it made of.}{}{fig:spectro}
If I could force the field to buy one instrument it would be this one, because a spectrometer
answers the question everything else dodges: what is it made of, or what is it burning.
Dispersing light into its component wavelengths yields emission and absorption features that
identify elements and molecules more or less uniquely. It is how we know the composition of
stars we will never visit. It is why a sodium street lamp, a mercury lamp, an LED, a magnesium
flare and an aurora cannot be confused by anybody holding a grating.

It also has the field's single best instrumented result, and it is not the result anybody
wanted. The Hessdalen project in Norway has monitored a valley in Holt\aa{}len continuously
since 1984, and Erling Strand's team, later joined by Massimo Teodorani, put spectrographs on
the lights. The optical spectra came back \emph{continuous} --- broadly thermal, without the
clean emission lines that would identify a specific ionised gas. That is a genuine, hard-won,
awkward finding. It constrains the plasma models that were the leading candidate, it excludes
several mundane light sources outright, and it delivers no verdict. Four decades of honest data
collection produced a narrowing, not an answer.

\pullquote{The Hessdalen programme's most valuable output is not a photograph or a spectrum. It
is the demonstration that a small group with modest funding can run a continuous instrumented
watch for forty years and publish what it finds --- including when what it finds is
inconvenient.}

A cheap diffraction grating over a phone camera will separate the common artificial light
sources and is worth carrying. A calibrated spectrograph needs a stable mount, a wavelength
calibration source, and somebody who understands that atmospheric absorption bands are in every
outdoor spectrum and are not a discovery.

\Needspace*{0.40\textheight}
\subsection{FLIR and Thermal Imaging}

\figwrap[r]{0.30}{A handheld thermal imaging camera. It measures received radiance, not temperature at range, and apparent size in the frame is dominated by blooming and focus.}{}{fig:flircam}
Forward-looking infrared sensors image thermal radiation, roughly 3--5 or 8--14 micrometres,
and they are the reason a great deal of modern UAP discussion exists at all. The 2004
\emph{Nimitz} and 2015 US Navy videos are ATFLIR recordings. Uncooled microbolometers are now
sold as phone accessories for a few hundred pounds.

Three things about thermal imagery that the field consistently gets wrong. First, \hi{a thermal
camera does not measure temperature at range}; it measures received radiance, and converting
that to a temperature requires knowing the emissivity of the target and the transmission of the
intervening atmosphere, neither of which you know for an unidentified object. Second, apparent
size in an infrared frame is dominated by sensor blooming and focus, so a small warm object and
a large cool one can look identical. Third, and most importantly, a gimballed targeting pod is a
moving optical assembly on a moving aircraft, and the artefacts it generates --- derotation
steps that look like an object spinning, glare structures that look like a hull, apparent
tangential velocity that is parallax against a wave-textured sea --- are not evidence of
anything except the mechanism producing them. I go through the specific cases in
Section~\ref{sec:uap}.

None of that makes the videos worthless, and I want to be careful here, because the debunking
has become as sloppy as the promotion. The derotation explanation accounts for the Gimbal
object's apparent rotation and does not obviously account for everything the operators reported
seeing, and ``sensor artefact'' offered without a specific mechanism is exactly the same
intellectual move as ``alien craft'' offered without one. The defensible position is that the
recordings are authentic sensor events, that the sensor is well enough understood to explain a
great deal of what looks strange in them, and that the remaining residual is small, real and
unresolved.

\Needspace*{0.40\textheight}
\subsection{UFODAP and Automated Multi-Sensor Stations}\label{sec:ufodapautomatedmultise}

\figwrap[l]{0.36}{An automated all-sky monitoring camera --- this one belongs to the FRIPON meteor network in Argentina. Unattended stations of this kind generate the quiet-hour denominator that expedition-based ufology has never had.}{}{fig:allsky}
The UFO Data Acquisition Project, developed principally by the electrical engineer and computer
scientist Ron Olch and deployed with Christopher O'Brien in Colorado's San Luis Valley, is the
most methodologically interesting hardware in civilian ufology. It is a multi-sensor station: a
pan/tilt/zoom camera with automated optical tracking, plus a sensor package typically including
a magnetometer, accelerometers, barometric pressure, temperature and humidity, an RF spectrum
monitor, and GPS for position and precise time. Detections trigger recording across all channels
simultaneously, with pre-trigger buffering, so you get the seconds \emph{before} the event as
well as after. Conceptually it is the descendant of the stalled UFODATA proposal, and it shares
a lineage with the Hessdalen automatic measurement station.

Why this design matters more than any individual sensor inside it: \hi{simultaneity across
independent physical channels is the only thing in this field that can defeat the standard
explanations.} A light in the sky is a light in the sky. A light in the sky with a coincident
magnetometer deflection, no correlated RF signature, timestamped to GPS, recorded automatically
by a station nobody was standing next to, is a data point that aircraft, lanterns, planets and
hoaxes struggle to produce. And critically, an automated station runs when nothing is happening,
which means it generates the denominator --- thousands of quiet hours --- that Section~\ref{sec:doubtdrivenskepticism}
insists on and that expedition-based ufology has never had.

The honest caveats. Cheap magnetometers drift with temperature and respond to passing vehicles,
mobile phones and the station's own electronics. Automated triggering produces enormous
false-positive volumes: insects, birds, satellites, aircraft strobes, lens contamination.
Unattended kit in remote terrain fails in ways nobody notices for weeks. And a network of three
stations would be worth far more than a dozen isolated ones, because triangulation is what turns
a bearing into a position and an angular rate into a speed. The instrumentation is no longer the
bottleneck. Funding, siting discipline and the willingness to publish the boring years are the
bottleneck.

\subsection{What Each Instrument Can and Cannot Settle}

\begin{center}\small
\begin{tabular}{@{}p{0.20\textwidth}p{0.33\textwidth}p{0.35\textwidth}@{}}
\textbf{Instrument} & \textbf{Can establish} & \textbf{Cannot establish}\\[3pt]
Geiger counter & Absence of elevated ionising radiation & The cause of any elevation; anything
at all without matched background\\[2pt]
GPR & Presence and geometry of subsurface disturbance & What caused the disturbance, or
when\\[2pt]
LIDAR & Range, surface geometry, a scattering track & Identity or composition; true velocity
from one station\\[2pt]
Spectrometer & Light-source type and composition; exclusion of known lamps and flares & Distance,
size or intent; anything uncalibrated through heavy atmosphere\\[2pt]
FLIR / thermal & That a radiance source was present, and its apparent motion in frame & Absolute
temperature, size, range or speed\\[2pt]
Multi-sensor station & Simultaneity across channels, precise timing, the quiet-hour denominator &
Position or velocity from a single station without triangulation\\
\end{tabular}
\end{center}

Read down the right-hand column, because that is the column the field skips. Every instrument
here is capable of producing a defensible negative result, and negative results are the
overwhelming majority of what honest instrumentation yields. If your equipment has never
returned a boring answer, either it is not being used properly or you are not publishing.

One last piece of advice, offered without irony. Buy the compass, the tape measure, the
calibrated grey card and the notebook before you buy anything with a digital display. A
photograph taken from a marked position with a known focal length, a bearing recorded and
corrected for declination, and a timestamp you can trust will resolve more cases than every
Geiger counter ever waved at a field in Suffolk.

And that is the Hynek ending. He is the man who set out with the best instrument of his era ---
an observatory, a professional reputation, and the confidence of the United States Air Force ---
and concluded that the useful equipment was a notebook, a sky reconstruction, and the willingness
to say what a report did and did not establish. \hi{Everything in this chapter is a footnote to
that conclusion.} The rest of the field is still buying displays.


\section{The Chemistry of Ufology}\label{sec:chemistryufology}

Almost every claim in this field that could in principle be settled tomorrow is a chemical
claim. Not a photograph, not a radar track, not a witness: a piece of matter, sitting in a
drawer, with a composition that can be measured. Chemistry is the only branch of physical
science the subject has ever handed a testable sample to, and it is therefore the place where
Ufology has come closest to behaving like a science and where it has failed most
embarrassingly. This section is about what a laboratory can actually establish from a fragment
of metal, what it cannot, and why the two best-known chemical claims in the literature ---
Element~115 and the ``impossible alloy'' --- collapse for opposite reasons.

Start with the instrument that makes the whole discussion possible. \hi{The periodic table is
not a catalogue of what has been found; it is a grid of predictions.} Its empty seats are
numbered, positioned and roughly described before anybody occupies them, which is why Dmitri
Mendeleev could leave gaps in 1869 and be proved right about gallium, scandium and germanium
within fifteen years. Everything an analytical chemist does to an anomalous sample is
ultimately a comparison against that grid, so it is worth having in front of you before the
claims start arriving.

% ---------------------------------------------------------------
% fig_ptable.tex --- the periodic table, for the Element 115 discussion in
% Section 4.6. GENERATED by tools/gen_ptable.py; edit the script, not this.
% Drawn in TikZ so it prints sharp at any size and carries no image credit.
% ---------------------------------------------------------------
\begingroup
\sffamily
\begin{center}
\refstepcounter{figph}\label{fig:ptable}%
% The whole picture is scaled down to the text width on the way in, so the
% sizes below are deliberately generous: after the squeeze the symbols land
% at about 7pt and the atomic numbers at about 4pt on the printed page.
\newcommand{\ptbox}[2]{\begin{tabular}{@{}c@{}}%
  {\fontsize{5.6}{5.6}\selectfont #1}\\[-0.4ex]%
  {\fontsize{9.6}{9.6}\selectfont\bfseries #2}\end{tabular}}
\resizebox{\linewidth}{!}{%
\begin{tikzpicture}[x=1cm,y=1cm,
  ptcell/.style={draw=midgrey,line width=0.3pt,minimum width=0.94cm,
                 minimum height=0.94cm,inner sep=0pt,text=ufogreenink},
  pt115/.style={draw=ufogreenink,line width=1.1pt,fill=ufogreen,
                minimum width=0.94cm,minimum height=0.94cm,inner sep=0pt,
                text=ufogreenink},
  legendswatch/.style={line width=0.4pt,minimum width=0.3cm,
                 minimum height=0.3cm,inner sep=0pt},
  legendtext/.style={anchor=west,text=ufogreenink,
                 font=\fontsize{8.4}{8.4}\selectfont}]
% group numbers along the top, period numbers down the left
\foreach \c in {1,...,18}{%
  \node[text=midgrey,font=\fontsize{6.2}{6.2}\selectfont] at (\c,-0.34) {\c};}
\foreach \r in {1,...,7}{%
  \node[text=midgrey,font=\fontsize{6.2}{6.2}\selectfont] at (0.3,-\r) {\r};}
% ---- the elements ----
\node[ptcell,fill=white] at (1,-1) {\ptbox{1}{H}};
\node[ptcell,fill=inkblue!14] at (18,-1) {\ptbox{2}{He}};
\node[ptcell,fill=amber!22] at (1,-2) {\ptbox{3}{Li}};
\node[ptcell,fill=amber!22] at (2,-2) {\ptbox{4}{Be}};
\node[ptcell,fill=ufogreendark!18] at (13,-2) {\ptbox{5}{B}};
\node[ptcell,fill=white] at (14,-2) {\ptbox{6}{C}};
\node[ptcell,fill=white] at (15,-2) {\ptbox{7}{N}};
\node[ptcell,fill=white] at (16,-2) {\ptbox{8}{O}};
\node[ptcell,fill=white] at (17,-2) {\ptbox{9}{F}};
\node[ptcell,fill=inkblue!14] at (18,-2) {\ptbox{10}{Ne}};
\node[ptcell,fill=amber!22] at (1,-3) {\ptbox{11}{Na}};
\node[ptcell,fill=amber!22] at (2,-3) {\ptbox{12}{Mg}};
\node[ptcell,fill=steel!8] at (13,-3) {\ptbox{13}{Al}};
\node[ptcell,fill=ufogreendark!18] at (14,-3) {\ptbox{14}{Si}};
\node[ptcell,fill=white] at (15,-3) {\ptbox{15}{P}};
\node[ptcell,fill=white] at (16,-3) {\ptbox{16}{S}};
\node[ptcell,fill=white] at (17,-3) {\ptbox{17}{Cl}};
\node[ptcell,fill=inkblue!14] at (18,-3) {\ptbox{18}{Ar}};
\node[ptcell,fill=amber!22] at (1,-4) {\ptbox{19}{K}};
\node[ptcell,fill=amber!22] at (2,-4) {\ptbox{20}{Ca}};
\node[ptcell,fill=steel!20] at (3,-4) {\ptbox{21}{Sc}};
\node[ptcell,fill=steel!20] at (4,-4) {\ptbox{22}{Ti}};
\node[ptcell,fill=steel!20] at (5,-4) {\ptbox{23}{V}};
\node[ptcell,fill=steel!20] at (6,-4) {\ptbox{24}{Cr}};
\node[ptcell,fill=steel!20] at (7,-4) {\ptbox{25}{Mn}};
\node[ptcell,fill=steel!20] at (8,-4) {\ptbox{26}{Fe}};
\node[ptcell,fill=steel!20] at (9,-4) {\ptbox{27}{Co}};
\node[ptcell,fill=steel!20] at (10,-4) {\ptbox{28}{Ni}};
\node[ptcell,fill=steel!20] at (11,-4) {\ptbox{29}{Cu}};
\node[ptcell,fill=steel!20] at (12,-4) {\ptbox{30}{Zn}};
\node[ptcell,fill=steel!8] at (13,-4) {\ptbox{31}{Ga}};
\node[ptcell,fill=ufogreendark!18] at (14,-4) {\ptbox{32}{Ge}};
\node[ptcell,fill=ufogreendark!18] at (15,-4) {\ptbox{33}{As}};
\node[ptcell,fill=white] at (16,-4) {\ptbox{34}{Se}};
\node[ptcell,fill=white] at (17,-4) {\ptbox{35}{Br}};
\node[ptcell,fill=inkblue!14] at (18,-4) {\ptbox{36}{Kr}};
\node[ptcell,fill=amber!22] at (1,-5) {\ptbox{37}{Rb}};
\node[ptcell,fill=amber!22] at (2,-5) {\ptbox{38}{Sr}};
\node[ptcell,fill=steel!20] at (3,-5) {\ptbox{39}{Y}};
\node[ptcell,fill=steel!20] at (4,-5) {\ptbox{40}{Zr}};
\node[ptcell,fill=steel!20] at (5,-5) {\ptbox{41}{Nb}};
\node[ptcell,fill=steel!20] at (6,-5) {\ptbox{42}{Mo}};
\node[ptcell,fill=steel!20,draw=gapred,dashed,line width=0.55pt] at (7,-5) {\ptbox{43}{Tc}};
\node[ptcell,fill=steel!20] at (8,-5) {\ptbox{44}{Ru}};
\node[ptcell,fill=steel!20] at (9,-5) {\ptbox{45}{Rh}};
\node[ptcell,fill=steel!20] at (10,-5) {\ptbox{46}{Pd}};
\node[ptcell,fill=steel!20] at (11,-5) {\ptbox{47}{Ag}};
\node[ptcell,fill=steel!20] at (12,-5) {\ptbox{48}{Cd}};
\node[ptcell,fill=steel!8] at (13,-5) {\ptbox{49}{In}};
\node[ptcell,fill=steel!8] at (14,-5) {\ptbox{50}{Sn}};
\node[ptcell,fill=ufogreendark!18] at (15,-5) {\ptbox{51}{Sb}};
\node[ptcell,fill=ufogreendark!18] at (16,-5) {\ptbox{52}{Te}};
\node[ptcell,fill=white] at (17,-5) {\ptbox{53}{I}};
\node[ptcell,fill=inkblue!14] at (18,-5) {\ptbox{54}{Xe}};
\node[ptcell,fill=amber!22] at (1,-6) {\ptbox{55}{Cs}};
\node[ptcell,fill=amber!22] at (2,-6) {\ptbox{56}{Ba}};
\node[ptcell,fill=midgrey!25] at (3,-9) {\ptbox{57}{La}};
\node[ptcell,fill=midgrey!25] at (4,-9) {\ptbox{58}{Ce}};
\node[ptcell,fill=midgrey!25] at (5,-9) {\ptbox{59}{Pr}};
\node[ptcell,fill=midgrey!25] at (6,-9) {\ptbox{60}{Nd}};
\node[ptcell,fill=midgrey!25,draw=gapred,dashed,line width=0.55pt] at (7,-9) {\ptbox{61}{Pm}};
\node[ptcell,fill=midgrey!25] at (8,-9) {\ptbox{62}{Sm}};
\node[ptcell,fill=midgrey!25] at (9,-9) {\ptbox{63}{Eu}};
\node[ptcell,fill=midgrey!25] at (10,-9) {\ptbox{64}{Gd}};
\node[ptcell,fill=midgrey!25] at (11,-9) {\ptbox{65}{Tb}};
\node[ptcell,fill=midgrey!25] at (12,-9) {\ptbox{66}{Dy}};
\node[ptcell,fill=midgrey!25] at (13,-9) {\ptbox{67}{Ho}};
\node[ptcell,fill=midgrey!25] at (14,-9) {\ptbox{68}{Er}};
\node[ptcell,fill=midgrey!25] at (15,-9) {\ptbox{69}{Tm}};
\node[ptcell,fill=midgrey!25] at (16,-9) {\ptbox{70}{Yb}};
\node[ptcell,fill=midgrey!25] at (17,-9) {\ptbox{71}{Lu}};
\node[ptcell,fill=steel!20] at (4,-6) {\ptbox{72}{Hf}};
\node[ptcell,fill=steel!20] at (5,-6) {\ptbox{73}{Ta}};
\node[ptcell,fill=steel!20] at (6,-6) {\ptbox{74}{W}};
\node[ptcell,fill=steel!20] at (7,-6) {\ptbox{75}{Re}};
\node[ptcell,fill=steel!20] at (8,-6) {\ptbox{76}{Os}};
\node[ptcell,fill=steel!20] at (9,-6) {\ptbox{77}{Ir}};
\node[ptcell,fill=steel!20] at (10,-6) {\ptbox{78}{Pt}};
\node[ptcell,fill=steel!20] at (11,-6) {\ptbox{79}{Au}};
\node[ptcell,fill=steel!20] at (12,-6) {\ptbox{80}{Hg}};
\node[ptcell,fill=steel!8] at (13,-6) {\ptbox{81}{Tl}};
\node[ptcell,fill=steel!8] at (14,-6) {\ptbox{82}{Pb}};
\node[ptcell,fill=steel!8] at (15,-6) {\ptbox{83}{Bi}};
\node[ptcell,fill=ufogreendark!18,draw=gapred,dashed,line width=0.55pt] at (16,-6) {\ptbox{84}{Po}};
\node[ptcell,fill=ufogreendark!18,draw=gapred,dashed,line width=0.55pt] at (17,-6) {\ptbox{85}{At}};
\node[ptcell,fill=inkblue!14,draw=gapred,dashed,line width=0.55pt] at (18,-6) {\ptbox{86}{Rn}};
\node[ptcell,fill=amber!22,draw=gapred,dashed,line width=0.55pt] at (1,-7) {\ptbox{87}{Fr}};
\node[ptcell,fill=amber!22,draw=gapred,dashed,line width=0.55pt] at (2,-7) {\ptbox{88}{Ra}};
\node[ptcell,fill=midgrey!25,draw=gapred,dashed,line width=0.55pt] at (3,-10) {\ptbox{89}{Ac}};
\node[ptcell,fill=midgrey!25,draw=gapred,dashed,line width=0.55pt] at (4,-10) {\ptbox{90}{Th}};
\node[ptcell,fill=midgrey!25,draw=gapred,dashed,line width=0.55pt] at (5,-10) {\ptbox{91}{Pa}};
\node[ptcell,fill=midgrey!25,draw=gapred,dashed,line width=0.55pt] at (6,-10) {\ptbox{92}{U}};
\node[ptcell,fill=midgrey!25,draw=gapred,dashed,line width=0.55pt] at (7,-10) {\ptbox{93}{Np}};
\node[ptcell,fill=midgrey!25,draw=gapred,dashed,line width=0.55pt] at (8,-10) {\ptbox{94}{Pu}};
\node[ptcell,fill=midgrey!25,draw=gapred,dashed,line width=0.55pt] at (9,-10) {\ptbox{95}{Am}};
\node[ptcell,fill=midgrey!25,draw=gapred,dashed,line width=0.55pt] at (10,-10) {\ptbox{96}{Cm}};
\node[ptcell,fill=midgrey!25,draw=gapred,dashed,line width=0.55pt] at (11,-10) {\ptbox{97}{Bk}};
\node[ptcell,fill=midgrey!25,draw=gapred,dashed,line width=0.55pt] at (12,-10) {\ptbox{98}{Cf}};
\node[ptcell,fill=midgrey!25,draw=gapred,dashed,line width=0.55pt] at (13,-10) {\ptbox{99}{Es}};
\node[ptcell,fill=midgrey!25,draw=gapred,dashed,line width=0.55pt] at (14,-10) {\ptbox{100}{Fm}};
\node[ptcell,fill=midgrey!25,draw=gapred,dashed,line width=0.55pt] at (15,-10) {\ptbox{101}{Md}};
\node[ptcell,fill=midgrey!25,draw=gapred,dashed,line width=0.55pt] at (16,-10) {\ptbox{102}{No}};
\node[ptcell,fill=midgrey!25,draw=gapred,dashed,line width=0.55pt] at (17,-10) {\ptbox{103}{Lr}};
\node[ptcell,fill=steel!20,draw=gapred,dashed,line width=0.55pt] at (4,-7) {\ptbox{104}{Rf}};
\node[ptcell,fill=steel!20,draw=gapred,dashed,line width=0.55pt] at (5,-7) {\ptbox{105}{Db}};
\node[ptcell,fill=steel!20,draw=gapred,dashed,line width=0.55pt] at (6,-7) {\ptbox{106}{Sg}};
\node[ptcell,fill=steel!20,draw=gapred,dashed,line width=0.55pt] at (7,-7) {\ptbox{107}{Bh}};
\node[ptcell,fill=steel!20,draw=gapred,dashed,line width=0.55pt] at (8,-7) {\ptbox{108}{Hs}};
\node[ptcell,fill=steel!20,draw=gapred,dashed,line width=0.55pt] at (9,-7) {\ptbox{109}{Mt}};
\node[ptcell,fill=steel!20,draw=gapred,dashed,line width=0.55pt] at (10,-7) {\ptbox{110}{Ds}};
\node[ptcell,fill=steel!20,draw=gapred,dashed,line width=0.55pt] at (11,-7) {\ptbox{111}{Rg}};
\node[ptcell,fill=steel!20,draw=gapred,dashed,line width=0.55pt] at (12,-7) {\ptbox{112}{Cn}};
\node[ptcell,fill=steel!8,draw=gapred,dashed,line width=0.55pt] at (13,-7) {\ptbox{113}{Nh}};
\node[ptcell,fill=steel!8,draw=gapred,dashed,line width=0.55pt] at (14,-7) {\ptbox{114}{Fl}};
\node[pt115] at (15,-7) {\ptbox{115}{Mc}};
\node[ptcell,fill=steel!8,draw=gapred,dashed,line width=0.55pt] at (16,-7) {\ptbox{116}{Lv}};
\node[ptcell,fill=ufogreendark!18,draw=gapred,dashed,line width=0.55pt] at (17,-7) {\ptbox{117}{Ts}};
\node[ptcell,fill=inkblue!14,draw=gapred,dashed,line width=0.55pt] at (18,-7) {\ptbox{118}{Og}};
% ---- the f-block, tied back to where it belongs ----
\node[draw=midgrey,line width=0.3pt,fill=midgrey!25,minimum width=0.94cm,
      minimum height=0.94cm,inner sep=0pt,text=midgrey,
      font=\fontsize{6.6}{6.6}\selectfont] at (3,-6) {57--71};
\node[draw=midgrey,line width=0.3pt,fill=midgrey!25,minimum width=0.94cm,
      minimum height=0.94cm,inner sep=0pt,text=midgrey,
      font=\fontsize{6.6}{6.6}\selectfont] at (3,-7) {89--103};
\node[text=midgrey,font=\fontsize{6.6}{7}\selectfont,anchor=east]
      at (2.45,-9) {lanthanides};
\node[text=midgrey,font=\fontsize{6.6}{7}\selectfont,anchor=east]
      at (2.45,-10) {actinides};
% ---- the Element 115 callout, set in the empty block above the d-block ----
\node[anchor=north west,text=ufogreenink,
      font=\fontsize{11.5}{12}\selectfont\bfseries] at (2.62,-0.58)
      {Element 115 --- moscovium (Mc)};
\node[anchor=north west,align=left,text width=9.5cm,text=ufogreenink,
      font=\fontsize{8.8}{10.6}\selectfont] at (2.62,-1.30)
      {Unmade and unobserved in 1989, when Lazar described a reactor running
       on it. Synthesised in 2003, named in 2016; its longest-lived known
       isotope lasts about two tenths of a second. The seat marked 115 was
       always going to be filled --- the table's blanks are numbered long
       before anyone reaches them.};
% ---- legend ----
\node[legendswatch,fill=amber!22,draw=midgrey] at (0.55,-10.85) {};
\node[legendtext] at (0.88,-10.85) {groups 1--2 metals};
\node[legendswatch,fill=steel!20,draw=midgrey] at (6.55,-10.85) {};
\node[legendtext] at (6.88,-10.85) {transition metals};
\node[legendswatch,fill=steel!8,draw=midgrey] at (12.55,-10.85) {};
\node[legendtext] at (12.88,-10.85) {other metals};
\node[legendswatch,fill=ufogreendark!18,draw=midgrey] at (0.55,-11.57) {};
\node[legendtext] at (0.88,-11.57) {metalloids};
\node[legendswatch,fill=white,draw=midgrey] at (6.55,-11.57) {};
\node[legendtext] at (6.88,-11.57) {non-metals and halogens};
\node[legendswatch,fill=inkblue!14,draw=midgrey] at (12.55,-11.57) {};
\node[legendtext] at (12.88,-11.57) {noble gases};
\node[legendswatch,fill=midgrey!25,draw=midgrey] at (0.55,-12.29) {};
\node[legendtext] at (0.88,-12.29) {lanthanides and actinides};
\node[legendswatch,fill=white,draw=gapred,dashed] at (6.55,-12.29) {};
\node[legendtext] at (6.88,-12.29) {no stable isotope};
\node[legendswatch,fill=ufogreen,draw=ufogreenink,line width=1pt] at (12.55,-12.29) {};
\node[legendtext] at (12.88,-12.29) {Element 115};
\end{tikzpicture}}\par\vspace{4pt}
{\footnotesize\color{steel}\textbf{\textcolor{ufogreenink}{Figure \thefigph.}} The
periodic table, with every element that has no stable isotope drawn in a broken
outline. Element 115 is the numbered seat in group 15, period 7 --- a blank the
table had reserved and described decades before anybody filled it.}
\end{center}
\endgroup


\subsection{Isotopes: the fingerprint that cannot be faked}

The single most useful fact in this entire book, for anyone who genuinely wants to know whether
a fragment came from somewhere else, is that the elements on that grid come in isotopes, and
that the proportions in which those isotopes occur are a fingerprint of origin.

An element is defined by its proton count; its isotopes differ in neutrons. Terrestrial
materials carry isotope ratios inherited from the formation of the solar nebula and modified
since by well-understood processes, and those ratios are astonishingly consistent. Ordinary
lead is about 52.4 per cent lead-208, 22.1 per cent lead-206, 24.1 per cent lead-207 and 1.4
per cent lead-204; magnesium is about 79 per cent magnesium-24, 10 per cent magnesium-25 and 11
per cent magnesium-26; titanium, silicon, chromium and oxygen all have equally fixed signatures.
A mass spectrometer measures these to parts per million, and a good laboratory can do it on a
milligram.

This matters because isotope ratios are the one property of matter that a hoaxer cannot adjust
and a terrestrial factory does not accidentally change. You can machine a strange shape, you can
sinter an unusual alloy, you can produce a layered composite in a garage with enough patience.
You cannot alter the isotopic composition of magnesium without an enrichment cascade, and the
facilities capable of separating isotopes at scale are counted in dozens worldwide and are all
known. \hi{An anomalous isotope ratio would be the strongest single piece of physical evidence
ever produced in the history of the subject.} It would not prove an intelligent origin --- a
presolar grain in a meteorite carries anomalous ratios and was made by a dying star, not a
spacecraft --- but it would establish, with no interpretive wriggle room whatsoever, that the
material did not come from Earth.

That is the test. It is cheap, it is available at any competent geochemistry department, it has
been available for sixty years, and the number of claimed UFO fragments that have been shown to
carry a non-terrestrial isotope signature is zero. The Ubatuba magnesium samples from Brazil,
tested repeatedly from the 1950s onwards and often described as unusually pure, were found to be
isotopically ordinary. The samples from Art Bell's ``Art's Parts'' affair, discussed in
Section~\ref{sec:artsparts}, were isotopically ordinary. The implant material recovered by
Roger Leir was isotopically ordinary. \hi{Every fragment that has ever reached a mass
spectrometer has come back terrestrial.} That is not a suppression of results; it is a result,
and it is the most important negative finding in the field.

\subsection{Impossible Bonds and Impossible Composites}

The second family of chemical claims is subtler and needs more care, because it is usually
phrased in a way that sounds like physics and is in fact a statement about manufacturing.

When a witness or an investigator says a material is impossible, they almost never mean that its
chemical bonds violate quantum mechanics. Bonding is one of the best-understood phenomena in
science: the electron configurations on the grid opposite determine what will bond with what, in
what ratio, with what geometry and at what energy, and no sample recovered from any UFO case has
ever contained a bond that the shell model forbids. What they mean, when the claim has any
content at all, is one of three much narrower things, and the three should never be confused.

The first is \emph{anomalous purity} --- a metal with fewer trace contaminants than an industrial
process of the claimed period could plausibly deliver. This is a real and measurable claim, and
it is also the weakest of the three, because purity is a moving target: zone refining and
float-zone silicon achieved figures in the 1950s that would have been called impossible in the
1930s, and a sample of unknown date proves nothing about the technology available when it was
made.

The second is an \emph{anomalous phase or microstructure} --- grain sizes, crystal orientations
or layer thicknesses that ordinary casting and rolling do not produce. This is a much better
class of claim, because microstructure records the thermal history of the material and thermal
histories are hard to fake convincingly. It is also where the field's most interesting single
sample sits: the layered bismuth and magnesium-zinc specimen of the ``Art's Parts'' affair, whose
alternating micron-scale bismuth films and magnesium-zinc layers were genuinely unusual when the
sample surfaced in 1996. The honest position, which I set out in full in
Section~\ref{sec:artsparts}, is that the structure is odd, that it has no established chain of
custody, and that comparable layered structures have since been produced by thin-film deposition
for thermoelectric research. An unusual microstructure with no provenance is a curiosity, not
evidence.

The third is a genuine \emph{immiscibility violation} --- two elements alloyed together that the
relevant phase diagram says will not mix. This is the strongest form of the claim and the one
worth taking seriously, because phase diagrams are experimental facts rather than opinions. It is
also the one that modern metallurgy has quietly dismantled. Mechanical alloying, rapid
solidification, ion implantation and additive manufacturing all routinely produce metastable
solid solutions of elements that will not mix in the melt. \hi{``The phase diagram says no'' is a
statement about equilibrium, and no manufactured object is at equilibrium.} A sample that
appeared impossible in 1955 may be ordinary laboratory practice in 2025, which is precisely why
the isotope test --- which does not depend on what any factory can or cannot do --- is the only
chemical argument in this field that does not decay with time.

The Element~115 story is the clearest illustration of the whole pattern, and I treat it in detail
alongside the Lazar account in Section~\ref{sec:s4lazar} and the reverse-engineering claim in
Section~\ref{sec:technologicalboom}. The short chemical version belongs here. Naming an
unoccupied seat in a numbered row is not a prediction; the grid had already numbered it. And the
grid also said what the occupant would be like: everything past bismuth is unstable, the
instability deepens with atomic number, and moscovium's longest-lived known isotope survives for
roughly two tenths of a second. \hi{An element that cannot survive long enough to be weighed
cannot fuel anything.} The half of the claim that came true is the half the table gave away for
free; the half that carried the entire story --- stability --- is the half chemistry ruled out in
advance.

\subsection{Metamaterials: Where Chemistry Stops Being the Answer}

There is one direction in which the ``impossible material'' intuition turns out to be pointing at
something real, and it is not chemistry at all. It is geometry, and it is the groundwork for
everything this book has to say later about propulsion, cloaking and the physics of the observed
flight characteristics.

A metamaterial is a structure whose behaviour comes from its architecture rather than its
composition. Build a lattice of metal rings and wires with a spacing much smaller than the
wavelength of the radiation you care about, and the assembly responds to that radiation as though
it were a homogeneous substance with properties no substance on the grid possesses --- including a
negative refractive index, which bends light the wrong way round. Victor Veselago worked out the
consequences theoretically in 1967. David Smith's group built one at microwave frequencies in
2000. Nothing in it is exotic: copper, fibreglass board, ordinary etching.

The reason this matters for Ufology is that it breaks the assumption on which every ``we would
need impossible matter to do that'' argument rests. \hi{A material's properties are not a
property of its elements alone.} Ask a chemist of 1950 what substance could have a negative
refractive index and the correct answer was none; ask what \emph{structure} could, and the answer
was a lattice you could have built in 1950 with the tools of 1950, had anyone known what to aim
for. Analyse a fragment of such a lattice by mass spectrometry and you will report copper and
epoxy. The composition tells you nothing. The information is in the arrangement, at a scale that
an elemental assay is blind to by construction.

This has two immediate consequences for how physical evidence should be handled, and both cut
against current practice in the field. First, destructive elemental analysis --- dissolving,
spark-testing or sectioning a sample to get a clean composition --- can destroy the only
information that mattered. \hi{The correct first pass on any anomalous fragment is
non-destructive and structural: micro-computed tomography, electron microscopy, X-ray
diffraction.} Composition second. The field has this order backwards almost universally, because
an elemental report is cheap and reads impressively.

Second, and more usefully, it tells you what an artefact of a substantially more advanced
technology would most plausibly look like on the bench: not a new element, which the grid says
cannot exist in a stable form anywhere in this universe, but a familiar element list in an
arrangement we cannot manufacture. Ordinary atoms, impossible architecture. That is a
falsifiable expectation, it is testable with instruments that exist in any materials department,
and as far as the public record goes it has never been tested on a UFO sample --- because nobody
looked, and because the samples with the most interesting structures are the ones with the worst
provenance.

I take up what such structures could actually do --- guiding fields, shaping radar returns,
managing the energy densities the manoeuvres in Chapter~\ref{ch:dimensions} appear to demand --- in
the chapters on propulsion and on the physics of the phenomenon. The point to carry forward is
the methodological one. Chemistry gives this field its one incorruptible test, in the isotopes,
and its one clear negative result, in the fact that every fragment ever measured has come back
terrestrial. Where chemistry runs out, the answer is not a stranger element. It is a stranger
arrangement of the elements we already have.


% ---------------------------------------------------------------
\section{The Biology of Ufology}\label{sec:biologyufology}

Chemistry hands this field one incorruptible test and one clear negative result. Biology hands it
two questions that are asked constantly and almost never asked properly, and both of them arrive
later in this book carrying more weight than the evidence usually gets to bear. The first is
whether anything not from here could interbreed with us, which is the load-bearing assumption of
every hybrid narrative from Genesis~6 to the modern abduction corpus. The second is why an object
lodged in living human tissue for years would sit there quietly instead of being attacked. Neither
question requires a spacecraft to be interesting, both are answerable at a bench, and the reason
they belong in a chapter about investigative method rather than in the chapters where the claims
appear is that a reader who does not already know the relevant biology cannot evaluate those claims
when they arrive. So the groundwork goes here.

\subsection{The 1.2 Per Cent That Cannot Be Crossed}

We share roughly 98.8 per cent of our DNA with chimpanzees. We diverged from them somewhere
between six and seven million years ago. We cannot produce a child with one, and the failure is
not marginal --- there is no recorded hybrid, viable or otherwise, despite a century of captive
proximity and at least one documented attempt.

That pairing is the single most useful calibration point in this entire discussion, because it
fixes the scale. \hi{A 1.2 per cent difference in sequence is already far too much.} Understanding
why requires separating three failures that the word ``compatible'' hides.

The first is architectural. Humans carry twenty-three chromosome pairs; chimpanzees, gorillas and
orangutans carry twenty-four. Our chromosome~2 is the fused product of two ancestral chromosomes,
and the fusion site is still visible in the sequence as an interstitial band of telomeric repeats
with a dead centromere beside it. Meiosis requires each chromosome to find and pair with its
partner along its length; a genome that has been rearranged relative to its counterpart cannot do
this cleanly, and the gametes that result carry the wrong complement of genetic material. This is
not a question of how similar the genes are. It is a filing problem, and it is fatal regardless of
content.

The second is molecular recognition at the moment of fertilisation. Sperm and egg identify each
other through specific protein pairs --- IZUMO1 on the sperm surface binding JUNO on the egg, and
the zona pellucida glycoproteins that surround the oocyte --- and these proteins are among the
fastest-evolving in the genome precisely because species-specific fertilisation is what keeps
species separate. Two lineages a few million years apart usually cannot achieve gamete binding at
all.

The third is developmental. Even where fertilisation succeeds, the embryo has to run a programme:
gene regulatory networks firing in sequence, imprinting patterns matching between the two parental
genomes, placental signalling negotiated between two organisms. The differences between us and
chimpanzees are concentrated in exactly this regulatory machinery rather than in the protein-coding
sequence, which is why the 98.8 per cent figure is so misleading to a lay reader. The parts list is
nearly identical. The assembly instructions are not.

Now put numbers on where the boundary actually sits, because it is not a mystery and it is not
arbitrary. Horses (2n = 64) and donkeys (2n = 62) diverged perhaps four million years ago; they
produce mules, which are robust and almost always sterile --- the chromosome mismatch permits
development but not meiosis. Lions and tigers, separated by something under four million years with
matching karyotypes, produce hybrids that are occasionally fertile in the female line, which is
Haldane's rule doing what it does. Our own species crossed with Neanderthals, from whom we had been
separated for perhaps five to seven hundred thousand years with an identical chromosome
architecture, and the offspring were fertile: every non-African population alive today carries
roughly one to two per cent Neanderthal DNA, and Denisovan sequence sits on top of that. So the
empirical picture is a staircase rather than a wall. Fertile crossing survives divergences measured
in hundreds of thousands of years and identical karyotypes. Sterile hybrids appear out to a few
million years with mismatched karyotypes. Past that, nothing, anywhere, in any taxon --- and we sit
comfortably past that from every other organism on Earth.

This is what makes the hybrid claim so much larger than the people making it usually realise. A
lineage with a genuinely independent origin would not merely be more different from us than a
chimpanzee is. It would fail at a level below any of the three mechanisms above, because
compatibility depends on details that are arbitrary rather than convergent: that the information is
carried in DNA at all, in right-handed double helices, using these four bases, read in triplets,
with this particular codon-to-amino-acid assignment out of the many that would work, packaged on
histones, in this number of chromosomes, with these fertilisation proteins. \hi{Convergent evolution
can produce an eye twice. It cannot produce the same genetic code twice, and reproduction requires
the code to match, not the anatomy.} Two organisms that look alike may be unable to exchange a
single functional gene. Two organisms that can produce a fertile child are, necessarily, relatives.

Which is why the argument that surrounds this subject is aimed in the wrong direction. Whether one
believes the accounts or not, the accounts themselves make a measurable claim, and the measurement
is available.

\actualq{whether interbreeding with non-human intelligences happened --- that is a question about
testimony, and testimony has been argued to a standstill for six decades without moving anybody.}{where
the threshold between reproductive compatibility and incompatibility actually lies, and what a
successful cross would therefore prove about shared ancestry. Viability is not a matter of degree.
It is a boundary with a location, and the location is knowable.}

The consequence is stated in full alongside Genesis~6 in Section~\ref{sec:nephilim} and alongside the
abduction hybrid material in Section~\ref{sec:geneticmotivation}, but the logic is short enough to
finish here. If a cross of that kind ever produced a viable child, the other party was not
biologically alien in any meaningful sense; it shared our ancestry, our chromosome architecture and
our genetic code, and the interesting problem becomes what we are rather than where they came from.
If the other party was genuinely of independent origin, the child is impossible and the accounts are
about something other than reproduction. Both branches are informative. Neither is available to
somebody who has not been told what the 1.2 per cent means.

\subsection{What the Immune System Actually Does to a Foreign Object}

The second piece of groundwork concerns the other physical claim in the abduction literature, and
it needs stating in advance for the same reason: the claim is remarkable only if you know what the
normal response looks like.

Human defence runs in two layers. The innate system is immediate and non-specific --- complement
proteins, neutrophils and macrophages responding to broad molecular patterns --- and the adaptive
system is slower, specific and remembers, using the major histocompatibility complex to present
fragments of anything it has digested to T cells that have been screened, during development, to
ignore self. That screening is why transplants fail. An organ from an unmatched human donor is
attacked because its MHC molecules are recognisably not ours, and this is worth registering before
any later claim about grafted or implanted living tissue, because non-human tissue would be a far
larger immunological event than the people describing it tend to assume. Hyperacute rejection of a
mismatched graft is measured in minutes.

A non-living object is handled differently, and the process has a name: the foreign body response.
Within seconds of insertion, plasma proteins adsorb onto the surface. Acute inflammation follows,
then macrophages arrive and attempt phagocytosis; finding the object too large to engulf, they fuse
into multinucleated foreign-body giant cells and secrete degradative enzymes against a surface they
cannot remove. Over the following weeks fibroblasts lay down an avascular collagen capsule,
typically tens to a few hundred micrometres thick, which walls the object off from the host.
\hi{Encapsulation is the default outcome for anything solid that ends up inside us, and the capsule
is a quarantine, not an acceptance.}

This is why an entire engineering discipline exists to negotiate with that reaction. No material is
truly inert; biocompatibility means provoking a manageable response rather than none. Titanium is
valuable in surgery not because the body ignores it but because its oxide surface allows bone to
grow into direct contact, and that integration took decades of deliberate work to characterise. A
splinter hurts, an unsterile fragment in soft tissue is normally a clinical problem, and an object
sitting for years in a human foot with no inflammatory infiltrate, no capsule of the expected kind
and a nerve-rich sheath adherent to its surface is --- if the observation is real --- a genuine
histological anomaly, whatever the object turns out to be made of.

That is precisely the observation the implant literature recorded and then discarded in favour of
arguing about metallurgy, and I take the failure apart in Section~\ref{sec:evidenceabductioncases}.
What belongs in a methods chapter is the protocol nobody ran. The tests are ordinary: blinded
histopathology by a pathologist who is not told what the specimen is supposed to be, capsule
thickness measured against matched controls of retained surgical debris, immunohistochemistry for
macrophages and T cells to establish whether an inflammatory population is present or absent, and
staining for nerve fibres to confirm or refute the claimed innervation. Every one of those is
routine in any hospital pathology department, none of them costs more than a few hundred dollars,
and as far as the public record shows not one of them has ever been performed on a claimed implant
under controlled conditions. \hi{The most interesting biological claim in Ufology has never been
given a biologist.}

\subsection{A Declaration of Competence}

I owe the reader a disclosure at the end of this, and it is the same disclosure I would want from
anybody handing me a chapter on a subject outside their training.

I am not a biologist. I did not know any of the above at the start of this project, I do not hold a
qualification that would let anyone take it on my authority, and none of it should be taken on my
authority. Before the arrival of systems capable of doing technical review at speed, a book like
this one had exactly two options at this point: leave the biology out, which is what almost every
book in the field has done, or hire a freelance biologist to check it, which costs money most
independent authors do not have and takes months they cannot spare. That is not a small structural
fact about the literature. It is one reason the hybrid question and the implant question have sat
unexamined for sixty years while the same photographs were argued about again and again.

So the standard I have held this section to is that nothing in it rests on me. Chromosome counts,
divergence estimates, the identity of the fertilisation proteins, the thickness of a fibrous
capsule, the timeline of the foreign body response --- these are textbook values with published
sources, listed at the end of this chapter, and any reader with a library card can check every one.
Where the underlying science is genuinely contested, such as the exact percentage identity between
the human and chimpanzee genomes, which shifts with the alignment method and whether insertions and
deletions are counted, I have said so rather than picking the number that sounds best. And if a
working biologist finds an error in this section, I want it: the correction goes into the next
edition with an acknowledgement, because being wrong in print and fixing it is how this is supposed
to work, and refusing to write about anything outside one's own expertise is how a field ends up
with nobody qualified to answer its most important question.

\begin{srcnotes}
\item Hynek, J.\ A.\ \emph{The UFO Experience: A Scientific Inquiry} (Henry Regnery, 1972):
the six-category observational classification, the close-encounter grades, the ``evidence about an
observation'' formulation and the residue argument used as this chapter's frame.
\item Hynek, J.\ A.\ \emph{The Hynek UFO Report} (Dell, 1977), for his retrospective account of
Category-A screening practice and the handling of the Trent (McMinnville) photographs.
\item Hynek, J.\ A.\ (21 October 1966). Letter to \emph{Science} on the scientific status of UFO
reports; and Center for UFO Studies (CUFOS) case-report series, Evanston, Illinois, 1974 onwards.
\item Mutual UFO Network, \emph{Field Investigator's Manual}, current edition, for the intake and
investigation procedures compared here against CUFOS practice.
\item Halt, C.\ I.\ (13 January 1981). \emph{Unexplained Lights}, memorandum, RAF Woodbridge;
and subsequent health-physics commentary on AN/PDR-27 sensitivity and the quoted 0.07 mR/hr
figure against 0.03--0.04 background.
\item Doleman, W.\ H.\ et al.\ \emph{The Roswell Dig Diaries} (2004), Sci-Fi Channel and the
University of New Mexico Office of Contract Archeology: remote sensing survey and excavation of
the claimed debris and skip sites on Bureau of Land Management land, former Foster Ranch.
\item Strand, E.\ \emph{Project Hessdalen} technical reports, 1984 onwards, and the Hessdalen
Automatic Measurement Station documentation.
\item Teodorani, M.\ (2004). A long-term scientific survey of the Hessdalen phenomenon.
\emph{Journal of Scientific Exploration}, 18(2); and Teodorani, M.\ \& Strand, E.\ Optical
spectrum analysis of the Hessdalen phenomenon (technical report).
\item UFODAP / UFO Data Acquisition Project technical and FAQ documentation, R.\ Olch; San Luis
Valley deployment with C.\ O'Brien.
\item Published analyses of the ATFLIR \emph{Nimitz} and GIMBAL recordings, together with the
opposing gimbal-derotation and glare-artefact analyses, for the range of positions summarised
here.
\item Isotope-ratio reference values: IUPAC Commission on Isotopic Abundances and Atomic
Weights, \emph{Atomic Weights of the Elements} and the isotopic-composition tables, current
release, for the terrestrial lead, magnesium, titanium and silicon signatures quoted.
\item Veselago, V.\ G.\ (1968). The electrodynamics of substances with simultaneously negative
values of permittivity and permeability. \emph{Soviet Physics Uspekhi}, 10(4); and Smith, D.\ R.\
et al.\ (2000). Composite medium with simultaneously negative permeability and permittivity.
\emph{Physical Review Letters}, 84(18).
\item Published analytical reports on the Ubatuba magnesium samples and on the layered
bismuth / magnesium-zinc ``Art's Parts'' specimen, together with the thin-film thermoelectric
deposition literature that produces comparable layer structures.
\item Oganessian, Yu.\ Ts.\ et al.\ (2004) on the synthesis of element 115, and IUPAC (2016) on
the naming of moscovium, for the isotope half-lives cited.
\item Human--chimpanzee genome comparison: Chimpanzee Sequencing and Analysis Consortium (2005).
Initial sequence of the chimpanzee genome and comparison with the human genome. \emph{Nature},
437(7055); and subsequent alignment-method critiques on how insertions and deletions alter the
quoted percentage identity.
\item Human chromosome~2 fusion: IJdo, J.\ W.\ et al.\ (1991). Origin of human chromosome 2: an
ancestral telomere--telomere fusion. \emph{Proceedings of the National Academy of Sciences},
88(20).
\item Gamete recognition proteins: Bianchi, E.\ et al.\ (2014). Juno is the egg Izumo receptor and
is essential for mammalian fertilization. \emph{Nature}, 508(7497); and the zona pellucida
glycoprotein literature on species-specific sperm binding.
\item Archaic introgression: Green, R.\ E.\ et al.\ (2010). A draft sequence of the Neandertal
genome. \emph{Science}, 328(5979); Reich, D.\ et al.\ (2010) on the Denisova hominin; and
subsequent estimates of one to two per cent Neanderthal ancestry in non-African populations.
\item Equid and felid hybrid sterility, karyotype counts and Haldane's rule, standard mammalian
cytogenetics and speciation references.
\item Foreign body response and biocompatibility: Anderson, J.\ M., Rodriguez, A.\ \& Chang, D.\ T.\
(2008). Foreign body reaction to biomaterials. \emph{Seminars in Immunology}, 20(2); and Ratner, B.\
D.\ et al.\ (eds), \emph{Biomaterials Science}, current edition, on protein adsorption,
foreign-body giant cells, capsule formation and titanium osseointegration.
\item Transplant immunology and MHC-mediated rejection: Janeway's \emph{Immunobiology}, current
edition, for the innate/adaptive division, self-tolerance and hyperacute rejection timelines cited.

\end{srcnotes}
